\section{Basic Probability: Thinking About Possibilities}

In addition to analyzing averages and relationships between variables,
statistics also involves probability. Probability helps us answer
questions such as:

\enquote{How likely is a particular condition to occur?}

In this study, there are \resN{} respondents, so these \resN{} respondents
can be treated as the sample space for this probability example.

To demonstrate the application of probability, the total X and Y scores
were grouped into a high category using a threshold of $\ge 60$.

Two events were defined:

A = a respondent has a high X score ($X \ge 60$)

B = a respondent has a high Y score ($Y \ge 60$)

Based on the dataset of \resN{} respondents:

\resHighX{} respondents had $X \ge 60$

\resHighY{} respondents had $Y \ge 60$

\resBoth{} respondents had both $X \ge 60$ and $Y \ge 60$

\subsection{Probability of A}

Formula:
\[
P(A) = \frac{\text{Number of respondents with high X}}{\text{Total number of respondents}}
\]
\[
P(A) = \frac{\resHighX}{\resN} = \resPARatio
\]
or approximately:
\[
P(A) = \resPA\%
\]
This means that, based on this dataset, the probability that a randomly
selected respondent has a high X score is approximately \resPA\%.

\subsection{Probability of B}

Formula:
\[
P(B) = \frac{\text{Number of respondents with high Y}}{\text{Total number of respondents}}
\]
\[
P(B) = \frac{\resHighY}{\resN} = \resPBRatio
\]
or approximately:
\[
P(B) = \resPB\%
\]
This means that the probability that a randomly selected respondent has a
high Y score is approximately \resPB\%.

\subsection{Probability of A and B}

We can also examine the probability that a respondent satisfies both
conditions simultaneously, meaning that the respondent has both a high X
score and a high Y score.

Formula:
\[
P(A \cap B) = \frac{\text{Number of respondents belonging to both A and B}}{\text{Total number of respondents}}
\]
\[
P(A \cap B) = \frac{\resBoth}{\resN} = \resPABRatio
\]
or approximately:
\[
P(A \cap B) = \resPAB\%
\]
This means that \resBoth{} out of \resN{} respondents had both X and Y
scores in the high category.

\subsection{Conditional Probability}

After calculating the probabilities of the two events, we can use
conditional probability to examine the probability of one event when
another condition is already known.

In this case, we want to know:

What is the probability that a respondent has a high Y score given that
the respondent has a high X score?

Formula:
\[
P(B \mid A) = \frac{P(A \cap B)}{P(A)}
\]
Using the research data:
\[
P(B \mid A) = \frac{\resBoth/\resN}{\resHighX/\resN}
\]
\[
P(B \mid A) = \frac{\resBoth}{\resHighX} = \resPBGivenARatio
\]
or approximately:
\[
P(B \mid A) = \resPBGivenA\%
\]
This means that among respondents with high X scores, approximately
\resPBGivenA\% also had high Y scores.

This example demonstrates how conditional probability uses additional
information to narrow the sample space. Before knowing the respondent's X
category, the probability of having a high Y score was \resPB\%. After
knowing that the respondent had a high X score, the conditional
probability of having a high Y score was \resPBGivenA\%.

Thus, the \resN{} respondents can be used not only to calculate averages
and
correlations, but also to demonstrate the concepts of probability,
intersection, and conditional probability using real research data.

\subsection{Probability Summary}

\begin{table}[htbp]
  \centering
  \singlespacing
  \caption{Probability Summary}
  \label{tab:probability}
  \begin{tabular}{ll}
    \toprule
    Calculation & Result \\
    \midrule
    Number of respondents & \resN \\
    High X ($\ge 60$) & \resHighX \\
    High Y ($\ge 60$) & \resHighY \\
    High X and High Y & \resBoth \\
    $P(A)$ & $\resHighX/\resN = \resPA\%$ \\
    $P(B)$ & $\resHighY/\resN = \resPB\%$ \\
    $P(A \cap B)$ & $\resBoth/\resN = \resPAB\%$ \\
    $P(B \mid A)$ & $\resBoth/\resHighX = \resPBGivenA\%$ \\
    \bottomrule
  \end{tabular}
\end{table}
